\scanpage{311}
\mspage{22\deleted{\unsure{6}}7}
% archive sheet number: 310

\diagram{\begin{tikzcd}[ampersand replacement=\&,column sep=large]
F_\alpha(j)\arrow[r,"u_\alpha(j)"]\arrow[dr,dashed,"p_\alpha(j)"']\&G(j)\arrow[r,no head,"v(j)"]\arrow[d,dashed,"\deleted{\unsure{p}}q(j)"]\&H(j)\arrow[dl,dashed,"q(j)"]\\
\&\unsure{\varphi|\phi}(j)\arrow[ul,phantom,"\text{comm.}"{near start}]\arrow[ur,phantom,"\text{comm.}"{near start}]\&
\end{tikzcd}}

(pour $j\in\inserted{J_0=}\mathrm{Ob}\,J$) provienne d'un \uline{foncteur}

\[w_\alpha:F_\alpha\longrightarrow H.\]

\deleted{C'est donc \unsure{en}} Pour prouver 4.9.3, nous sommes donc ramenés à prouver les deux \unsure{propositions|propriétés} $\mathrm{F}1$), $\mathrm{F}3_{\unsure{\pi|\kappa}})$ \unsure{précédentes}, sous l'une des deux conditions a) ou b) de 4.9.3.

\deleted{4.9.6.5. \uline{Démonstration de} 4.9.6.3}

\uline{Corollaire} 4.9.6.5 (de 4.9.6.3). Soit $(f_\alpha:I_\alpha\unsure{\mathbin{-}|\longrightarrow} I'_\alpha)_{\alpha\in\unsure{A}}$ une famille de foncteurs cofinaux, les \deleted{\unsure{catégories}}$I_\alpha$ sont \deleted{\unsure{étant}} filtrantes. Alors le foncteur produit $f=\prod f_\alpha:I=\prod_\alpha I_\alpha\longrightarrow I'=\prod_\alpha I'_\alpha$ est cofinal.

En effet, \deleted{\unsure{On sait que $I$ est filtrante|On sait que ces foncteurs sont},} \inserted{il faut vérifier que les} \deleted{catégories} $ {}_{i'}\backslash I$, pour $i'\in\mathrm{Ob}\,I'$, sont $0$-connexes. Or si $i'=(i'_\alpha)_{\alpha\in A}$, on a

\[ {}_{i'}\backslash I\simeq\prod_\alpha({}_{i'_\alpha}\backslash I_\alpha),\]

et par 4.9.6.3, 1\textsuperscript{o}) $\unsure{a})\Rightarrow c)$, les $ {}_{i'_\alpha}\backslash I_\alpha$ sont filtrantes, donc leur produit est filtrant

